\subsection{A short history of agents and MAS}
\label{sec:agents-history}
\slidesrc{Intro to agents \& MAS: 42--66}

Current interest in autonomous agents did not emerge from a vacuum. The main contributing
disciplines are \textbf{artificial intelligence}, \textbf{object-oriented programming} (and
concurrent object-based systems) and \textbf{human-computer interface design} (slide 43).

\subsubsection{From artificial intelligence}
\slidesrc{Intro to agents \& MAS: 44--47}

\begin{center}
\begin{tikzpicture}[font=\small, >=Stealth,
  ev/.style={anchor=north west, text width=10.2cm, align=left, inner sep=2pt},
  yr/.style={anchor=north east, font=\small\bfseries, inner sep=2pt}]
  \coordinate (e0) at (0.3,0);
  \foreach \i/\year/\txt/\c in {
    1/1961--71/{AI \textbf{planning}: Newell and Simon (1961), Fikes and Nilsson's STRIPS (1971). Until the 1980s, AI studies the components of intelligence (learning, reasoning, vision\ldots) \emph{in isolation}; expert systems are ``disembodied intelligence'': no real-time capabilities, no social ability.}/nbblue,
    2/1987/{Disillusion with \textbf{first-principle planning} (Chapman): building a new plan from scratch for every goal is too slow to be reactive (``calculative rationality'') and undecidable in many cases. Georgeff and Lansky propose the \textbf{BDI} architecture for practical reasoning.}/nbred,
    3/1991/{Brooks: \textbf{behavioural} or \textbf{reactive AI}, intelligence without symbolic representation.}/nborange,
    4/1992--94/{\textbf{Hybrid architectures}, usually stratified in layers: TouringMachines (Ferguson, 1992), InteRRaP (M\"uller and Pischel, 1994).}/nbpurple,
    5/1993/{Shoham, \emph{Agent-oriented programming}: agents are not just a branch of AI, but an autonomous technological and methodological approach. Shoham is considered the father of agent-oriented programming.}/nbgreen}
  {
    \pgfmathtruncatemacro{\prev}{\i-1}
    \node[ev] (e\i) at ([yshift=-4pt]e\prev.south west) {\txt};
    \node[yr, text=\c] at ([xshift=-0.55cm]e\i.north west) {\year};
    \fill[\c] ([xshift=-0.3cm, yshift=-7pt]e\i.north west) circle (3.2pt);
  }
  \begin{scope}[on background layer]
    \draw[very thick, nbblue!50] (0,0) -- (0,0 |- e5.south);
  \end{scope}
\end{tikzpicture}
\captionof{figure}{From AI to agents (slides 44--47). The architectures named here (reactive,
deliberative, hybrid, BDI) are the subject of Section~\ref{sec:agent-architectures}.}
\end{center}

\begin{nbdef}{Practical reasoning agents (slide 46)}
Agents whose architecture is modelled on, or inspired by, a theory of \textbf{practical
reasoning} in humans: the kind of pragmatic reasoning that we use to decide \emph{what to
do} (as opposed to theoretical reasoning, which decides what is true). The most widely known
practical reasoning architecture is \textbf{Belief-Desire-Intention} (BDI).
\end{nbdef}

\subsubsection{From object-oriented programming}
\slidesrc{Intro to agents \& MAS: 48--49}

Objects are computational entities that \emph{encapsulate some state}, are able to
\emph{perform actions} (methods) on this state, and \emph{communicate by message passing}.
This sounds very much like an agent. So, what is the difference?

\begin{center}
\small
\begin{tabularx}{\linewidth}{@{}l >{\raggedright\arraybackslash}X >{\raggedright\arraybackslash}X@{}}
\toprule
 & \textbf{Object} & \textbf{Agent} \\
\midrule
Autonomy & controls its \emph{state} (private fields), but not its \emph{behaviour}: if
  another object calls a public method, the method runs & controls its state \emph{and} its
  behaviour: a request is a message the agent may refuse \\
Motto & \multicolumn{2}{c}{``Objects do it for free; agents do it for money.''} \\
Flexibility & the standard object model says nothing about how to build systems that
  integrate proactive, reactive and social behaviour & reactive, proactive and social
  behaviour are part of the model \\
Thread of control & usually passive: acts only when called & has its own thread of
  control: always running, takes the initiative \\
\bottomrule
\end{tabularx}
\captionof{table}{Agents vs.\ objects (slide 49; last row added).}
\end{center}

\begin{nbexample}{the same request, two worlds}
\notinslides
Object: \code{account.withdraw(100)} \emph{always} executes the method (perhaps throwing
an exception). Agent: Alice's banking agent \emph{receives} the message ``please transfer
100 euro to Bob'', checks its own goals (``keep at least 500 on the account'', ``do not pay
unknown parties'') and may answer ``refuse''. One could argue that we can program this with
objects too; the point (slide 49) is that the object model has nothing to say about it,
while for agents it is the core of the abstraction.
\end{nbexample}

\subsubsection{From human-computer interaction}
\slidesrc{Intro to agents \& MAS: 50}

Instead of \emph{direct manipulation} (the user does everything by clicking), we would like
programs that behave like \textbf{``digital butlers''} (N.~Negroponte, \emph{Being
Digital}, 1995):
\begin{itemize}[nosep]
  \item ``The agent answers the phone, recognizes the callers, disturbs you when
        appropriate, and may even tell a white lie on your behalf.''
  \item ``If you have somebody who knows you well and shares much of your information,
        that person can act on your behalf very effectively. If your secretary falls ill, it
        would make no difference if the temping agency could send you Albert Einstein. This
        issue is not about IQ. It is shared knowledge and the practice of using it in your
        best interests.''
  \item ``Like an army commander sending a scout ahead \ldots\ you will dispatch agents to
        collect information on your behalf.''
\end{itemize}
Thirty years later, personal assistants and AI agents that read your e-mail and book your
travel are trying to fulfil exactly this vision.

\subsubsection{Agents as intentional systems}
\slidesrc{Intro to agents \& MAS: 51--60}

When explaining human activity we often say things like:
\begin{itemize}[nosep]
  \item Janine took her umbrella because she \emph{believed} it was going to rain;
  \item Michael worked hard because he \emph{wanted} to possess a PhD.
\end{itemize}
These statements use a \textbf{folk psychology}: human behaviour is predicted and explained
by attributing attitudes such as believing, wanting, hoping, fearing. These attitudes are
called the \textbf{intentional notions}.

\begin{nbdef}{Intentional system (D.~C.~Dennett, \emph{The Intentional Stance}, 1989)}
An \textbf{intentional system} is a system made up of entities whose behaviour can be
predicted by the method of attributing \emph{beliefs}, \emph{desires} and \emph{rational
acumen}.
\end{nbdef}

\begin{nbintuition}[The intentional stance in practice]
You do not know the code of a chess program, but you predict its moves by saying ``it
\emph{wants} to win and \emph{believes} my queen is unprotected, so it will take it''. That
works much better than reasoning about transistors. Taking the intentional stance towards a
system is useful when it is too complex to explain mechanically, yet behaves rationally.
\end{nbintuition}

If we adhere to the \emph{strong} definition of agents, intentional systems are a suitable
theory for agents (slides 56--60). Representing intentional notions raises delicate
syntactic and semantic questions, and the usual answers are:
\begin{itemize}[nosep]
  \item \textbf{modal logic} languages are suitable for \emph{specifying} agents as
        intentional systems;
  \item languages based on \textbf{computational (modal) logic} are suitable for
        \emph{programming} agents;
  \item axiomatised logic-based languages can undergo \textbf{axiomatic verification};
        other languages can be used for \textbf{model checking}.
\end{itemize}

\begin{nbdef}{Modal logic (slides 61--63)}
Modal logic extends classical logic with a new connective $\Box$ (\textbf{necessity}) and
its derivable counterpart $\Diamond$ (\textbf{possibility}).
\begin{itemize}[nosep]
  \item $\Box p$: $p$ is \emph{necessarily} true, i.e.\ true in every possible scenario
        (``possible world'');
  \item $\Diamond p$: $p$ is \emph{possibly} true, i.e.\ true in at least one possible
        scenario.
\end{itemize}
$\Diamond$ can be defined from $\Box$: \quad $\Diamond p \liff \neg\Box\neg p$.
\end{nbdef}

``$p$ is possible'' = ``it is not necessary that $p$ is false''. It is the same duality as
$\exists$ and $\forall$ (Section~\ref{sec:quantifiers}): $\Box$ is a ``for all worlds'',
$\Diamond$ a ``there exists a world''.

\begin{center}
\small
\begin{tabularx}{\linewidth}{@{}l l >{\raggedright\arraybackslash}X@{}}
\toprule
\textbf{Modal logic} & \textbf{$\Box p$ reads} & \textbf{Example} \\
\midrule
epistemic & the agent knows/believes $p$ & $B_{\mathit{robot}}\,\mathit{door\_open}$: the robot believes the door is open \\
temporal & $p$ always holds (in the future) & $\Box\neg\mathit{collision}$: the drones never collide \\
deontic & $p$ is obligatory & $O(\mathit{pay\_taxes})$; $\Diamond$ becomes ``permitted'' \\
dynamic & after every execution of action $a$, $p$ holds & $[\mathit{open}]\,\mathit{door\_open}$ \\
\bottomrule
\end{tabularx}
\captionof{table}{Kinds of modal logics (slide 64; readings and examples added). Combinations
exist: \textbf{BDI logic} combines beliefs, desires, intentions and time; \textbf{KARO}
combines knowledge, abilities, results and opportunities.}
\end{center}

\subsubsection{Two ancestors of MAS: Contract Net and Actors}
\slidesrc{Intro to agents \& MAS: 65--66}

\begin{nbdef}{Contract Net Protocol (Davis and Smith, 1983)}
Agents can dynamically take two \textbf{roles}: \textbf{manager} or \textbf{contractor}.
\begin{enumerate}[nosep]
  \item Given a task, an agent first determines whether it can break it into subtasks that
        could be performed concurrently.
  \item As \emph{manager}, it uses the protocol to \textbf{announce} the tasks that could
        be transferred, and requests \textbf{bids} from nodes that could perform them.
  \item A node that receives an announcement replies with a \textbf{bid}, indicating how
        well it thinks it can perform the task.
  \item The manager collects the bids and \textbf{awards} the task to the best bidder,
        which becomes a \emph{contractor}.
\end{enumerate}
\end{nbdef}

\begin{center}
\begin{tikzpicture}[font=\small, >=Stealth, x=1.05cm, y=-0.8cm]
  \foreach \x/\n in {0/Manager, 3.6/Contractor A, 6.6/Contractor B, 9.6/Contractor C}{
    \node[draw, rounded corners, thick, fill=nbblue!8, minimum width=2.2cm] at (\x,0) {\n};
    \draw[dashed, nbgray] (\x,0.6) -- (\x,8.6);
  }
  \draw[->, thick, nbblue] (0,1.3) -- node[above, pos=0.2, font=\scriptsize]{announce task $T$} (3.6,1.3);
  \draw[->, thick, nbblue] (0,1.3) -- (6.6,1.6);
  \draw[->, thick, nbblue] (0,1.3) -- (9.6,1.9);
  \draw[->, thick, nborange] (3.6,3.2) -- node[above, font=\scriptsize]{bid: 5 min} (0,3.2);
  \draw[->, thick, nborange] (6.6,4.0) -- node[above, pos=0.7, font=\scriptsize]{bid: 2 min} (0,4.0);
  \draw[->, thick, nborange] (9.6,4.8) -- node[above, pos=0.8, font=\scriptsize]{bid: 9 min} (0,4.8);
  \node[draw, fill=nbgreen!10, font=\scriptsize, align=center] at (0,6.0) {evaluate\\bids};
  \draw[->, thick, nbgreen!70!black] (0,7.2) -- node[above, pos=0.6, font=\scriptsize]{award $T$} (6.6,7.2);
  \draw[->, thick, nbgray] (6.6,8.2) -- node[above, pos=0.5, font=\scriptsize]{result} (0,8.2);
\end{tikzpicture}
\captionof{figure}{A run of the Contract Net Protocol: one announcement, three bids, the task is
awarded to B (the best bid). The same protocol appears in FIPA as an interaction protocol
(Section~\ref{sec:agent-communication}).}
\end{center}

\begin{nbwarning}[Slide 65: who collects the bids?]
The slide says ``\emph{the contractor} collects the bids and awards the task to the best
bidder''. It is the \textbf{manager} who announces the task, collects the bids and awards
the task; the winner then becomes the contractor.
\end{nbwarning}

\begin{nbwhy}
Contract Net is everywhere where work is distributed by auction: a cloud scheduler choosing
the server that can run a job fastest, a logistics platform assigning a delivery to the
nearest courier, a factory controller assigning a part to the least loaded machine.
\end{nbwhy}

\begin{nbdef}{Actors (Hewitt, Bishop, Steiger, IJCAI 1973; Agha, 1986)}
Actors were proposed as universal primitives of concurrent computation: self-contained,
interactive, autonomous components of a computing system that communicate by
\textbf{asynchronous message passing}. The basic actor primitives are:
\begin{itemize}[nosep]
  \item \textbf{create}: create an actor from a behaviour description and a set of
        parameters (possibly including existing actors);
  \item \textbf{send}: send a message to an actor;
  \item \textbf{become}: change an actor's local state.
\end{itemize}
\end{nbdef}

\begin{nbintuition}[Actors today]
Each actor has a mailbox, processes one message at a time, and can only affect others by
sending messages: no shared memory, so no locks. Erlang (used for telecom switches and
WhatsApp's servers) and the Akka framework are built on actors. Agents add to actors
\emph{goals} and \emph{reasoning}; the messaging core is the same, and JADE agents
(Section~\ref{sec:jade}) work exactly like this.
\end{nbintuition}
